% ======================================================================
% Chapter 15: Non-Abelian Gauge Invariance
% (book pages 481-503; PDF pages 504-527)
% ======================================================================
\chapter{非阿贝尔规范不变性}

迄今为止，在本书中我们只处理了一类相当有限的量子场与相互作用，把注意力局限于标量场论、Yukawa 理论和量子电动力学。毫不奇怪，这些理论不足以描述基本粒子的全部已知相互作用。但是，既然可重整化理论的拉格朗日量不能包含质量维数高于 4 的项，那么还可能有其他什么理论呢？

接下来最自然的尝试是那些具有矢量场之间相互作用的理论，其形式为 $A^2(\partial A)$ 或 $A^4$。然而，这类合理的理论难以构造，因为矢量场算符的时间分量 $A^0$ 会给出负模方态。在第~5.5 节我们看到了这些负模方态在 QED 中并不造成困难：凭借 Ward 恒等式，它们被纵向极化态有效地抵消。而 Ward 恒等式又源于 QED 拉格朗日量在局域规范变换下的不变性。因此，也许如果我们能够推广局域规范不变性的原理，它就会引导我们构造出矢量粒子理论的其他合理形式。

本章的目标正是做这件事。首先我们将简短地回到对 QED 的研究，这一次把规范对称性视为根本原理，并从这一原理导出其余的理论。然后，在第~15.2 节中我们将看到，电动力学的规范不变性只是一个无穷参数对称性最平凡的例子，而更一般的例子会导出其他有趣的拉格朗日量。这些场论以深刻的方式推广了电动力学，其中第一个是由 Yang 和 Mills 构造的，%
\footnote{C.~N.~Yang and R.~Mills, \emph{Phys.\ Rev.} \textbf{96}, 191 (1954)。}
它们是多个矢量粒子的理论，其相互作用受到对称性原理的强烈约束。在后续章节中，我们将研究这些理论的量子化及其在基本粒子物理现实世界中的应用。

\section{规范不变性的几何}

在第~4.1 节中我们写下了量子电动力学的拉格朗日量，并注意到一个奇特的事实：它在变换 \eqref{eq:4.6} 的一个非常大的群下不变，允许在时空的每一点做独立的对称变换。这一不变性就是 QED 著名的\emph{规范对称性}。然而，从现代的观点看，规范对称性并不是一个偶然的趣事，而是决定拉格朗日量形式的根本原理。现在让我们以现代观点来回顾这个理论的各个要素。

我们从复值 Dirac 场 $\psi(x)$ 出发，规定我们的理论应在变换
\begin{equation}
  \psi(x) \to e^{i\alpha(x)}\psi(x) .
  \label{eq:15.1}
\end{equation}
下不变。这是一个转动角 $\alpha(x)$ 逐点任意变化的相位转动。我们怎样才能写出在此变换下不变的拉格朗日量呢？只要考虑拉格朗日量中不含导数的项，这很容易：我们只需写出那些对整体相位转动不变的同样的项。例如，费米子质量项
\begin{equation}
  m\bar{\psi}(x)\psi(x)
\end{equation}
被整体相位不变性所允许，而局域不变性不再给出进一步限制。

当我们试图写出包含导数的项时，困难就出现了。$\psi(x)$ 沿矢量 $n^\mu$ 方向的导数由极限过程
\begin{equation}
  n^\mu\partial_\mu\psi(x) = \lim_{\epsilon\to 0}\ \frac{1}{\epsilon}\bigl[
    \psi(x+\epsilon n) - \psi(x) \bigr] .
  \label{eq:15.2}
\end{equation}
定义。然而，在具有局域相位不变性的理论中，这个定义并不很合理，因为被相减的两个场 $\psi(x+\epsilon n)$ 和 $\psi(x)$ 在对称性 \eqref{eq:15.1} 下具有完全不同的变换。换言之，量 $\partial_\mu\psi$ 没有简单的变换律，也没有有用的几何解释。

为了以有意义的方式相减 $\psi(x)$ 在相邻点的取值，我们必须引入一个因子，它补偿从一点到下一点相位变换的差异。做到这一点的最简单方式是定义一个依赖于这两个点的标量 $U(y,x)$，它的变换律为
\begin{equation}
  U(y,x) \to e^{i\alpha(y)}U(y,x)e^{-i\alpha(x)}
  \label{eq:15.3}
\end{equation}
与 \eqref{eq:15.1} 同时成立。在零间隔处，我们设 $U(y,y) = 1$；一般而言，我们可以要求 $U(y,x)$ 是一个纯相位：$U(y,x) = \exp[i\phi(y,x)]$。有了这个定义，$\psi(y)$ 与 $U(y,x)\psi(x)$ 就具有相同的变换律，我们就可以以在局域对称性下仍然有意义的方式将它们相减。于是我们可以定义一个合理的导数，称为\emph{协变导数}，如下：
\begin{equation}
  n^\mu D_\mu\psi(x) = \lim_{\epsilon\to 0}\ \frac{1}{\epsilon}\bigl[
    \psi(x+\epsilon n) - U(x+\epsilon n, x)\psi(x) \bigr] .
  \label{eq:15.4}
\end{equation}
为了使这个定义明确，我们需要比较器 $U(y,x)$ 在无穷小间隔两点处的表达式。如果 $U(y,x)$ 的相位是位置 $y$ 和 $x$ 的连续函数，那么 $U(y,x)$ 可以在两点的间隔上展开：
\begin{equation}
  U(x+\epsilon n, x) = 1 - ie\epsilon n^\mu A_\mu(x) + \mathcal{O}(\epsilon^2) .
  \label{eq:15.5}
\end{equation}
这里我们任意地提取了一个常数 $e$。位移 $\epsilon n^\mu$ 的系数是一个新的矢量场 $A_\mu(x)$。这样的场作为局域对称性变换之比较器的无穷小极限而出现，称为\emph{联络}。于是协变导数取如下形式：
\begin{equation}
  D_\mu\psi(x) = \partial_\mu\psi(x) + ieA_\mu(x)\psi(x) .
  \label{eq:15.6}
\end{equation}
把 \eqref{eq:15.5} 代入 \eqref{eq:15.3}，人们发现 $A_\mu$ 在此局域规范变换下按如下方式变换：
\begin{equation}
  A_\mu(x) \to A_\mu(x) - \frac{1}{e}\partial_\mu\alpha(x) .
  \label{eq:15.7}
\end{equation}
为了检验所有这些表达式是自洽的，我们可以按照式~\eqref{eq:15.1} 和 \eqref{eq:15.7} 变换 $D_\mu\psi(x)$：
\begin{equation}
  D_\mu\psi(x) \to \Bigl( \partial_\mu + ie\Bigl( A_\mu - \frac{1}{e}
    \partial_\mu\alpha \Bigr) \Bigr)\ e^{i\alpha(x)}\psi(x)
  = e^{i\alpha(x)}\bigl( \partial_\mu + ieA_\mu(x) \bigr)\psi(x)
  = e^{i\alpha(x)}D_\mu\psi(x) .
  \label{eq:15.8}
\end{equation}
因此协变导数与场 $\psi$ 本身按相同的方式变换，正如我们构造它时在原初定义 \eqref{eq:15.4} 中所要求的那样。

我们现在已经重新得到了 QED 拉格朗日量的大部分熟悉成分。然而，从我们当前的观点看，协变导数的定义和联络 $A_\mu$ 的变换律都来自局域相位转动对称性的假设。甚至矢量场 $A_\mu$ 的存在本身也是局域对称性的结果：没有它，我们就不可能写出涉及 $\psi$ 之导数的不变拉格朗日量。

更一般地，我们当前的分析给出了一种构造所有在局域对称性下不变的拉格朗日量的方法。在任何含 $\psi$ 之导数的项中，把这些导数替换为协变导数。根据式~\eqref{eq:15.8}，它们与 $\psi$ 本身严格按相同的方式变换。因此，任何在整体相位转动下不变的 $\psi$ 及其协变导数的组合（而且只有这些组合）也将是局域不变的。

为了完成局域不变拉格朗日量的构造，我们必须找到场 $A_\mu$ 的动能项：一个依赖 $A_\mu$ 及其导数但不依赖 $\psi$ 的局域不变项。这一项既可以整体地由比较器 $U(y,x)$ 构造，也可以无穷小地由协变导数构造。

\begin{figure}[t]
  \centering
  \includegraphics[width=0.5\textwidth]{fig15.1.png}
  \caption{通过在 (1,2) 平面内围绕一个小方块的比较来构造场强。}
  \label{fig:15.1}
\end{figure}

从 $U(y,x)$ 出发，我们将需要把显式公式 \eqref{eq:15.5} 推广到 $\epsilon$ 展开式的下一项。利用 $U(y,x)$ 是纯相位的假设和限制条件 $U(x,y)^{\dagger} = U(y,x)$，可得：
\begin{equation}
  U(x+\epsilon n, x) = \exp\Bigl[ -ie\epsilon n^\mu A_\mu\Bigl( x +
    \tfrac{1}{2}\epsilon n \Bigr) + \mathcal{O}(\epsilon^3) \Bigr] .
  \label{eq:15.9}
\end{equation}
（放宽这些限制会在理论中引入额外的矢量场；这是一种不必要的复杂化。）利用 $U(y,x)$ 的这个展开式，我们把围绕时空中的一个小方块的相位方向比较链接起来。为明确起见，我们取这个小方块位于由单位矢量 $\hat{1}$、$\hat{2}$ 定义的 $(1,2)$ 平面内（见图~\ref{fig:15.1}）。定义 $U(x)$ 为围绕该圈四个角上的四次比较之乘积：
\begin{equation}
  U(x) = U(x, x+\epsilon\hat{2})\ U(x+\epsilon\hat{2}, x+\epsilon\hat{1}
    +\epsilon\hat{2})\ U(x+\epsilon\hat{1}+\epsilon\hat{2},
      x+\epsilon\hat{1})\ U(x+\epsilon\hat{1}, x) .
  \label{eq:15.10}
\end{equation}
$U$ 的变换律 \eqref{eq:15.3} 意味着 $U(x)$ 是局域不变的。因此在 $\epsilon\to 0$ 的极限下，它将给出 $A_\mu$ 的一个局域不变函数。为了求出这个函数的形式，代入展开式 \eqref{eq:15.9} 得到：
\begin{equation}
  U(x) = \exp\Bigl\{ -ie\epsilon\Bigl[ -A_2\Bigl( x+\tfrac{\epsilon}{2}\hat{2}
    \Bigr) - A_1\Bigl( x + \epsilon\hat{1} + \tfrac{\epsilon}{2}\hat{2} \Bigr)
      + A_2\Bigl( x + \epsilon\hat{1} + \tfrac{\epsilon}{2}\hat{2} \Bigr)
        + A_1\Bigl( x + \tfrac{\epsilon}{2}\hat{1} \Bigr) \Bigr]
          + \mathcal{O}(\epsilon^4) \Bigr\} .
  \label{eq:15.11}
\end{equation}
当我们把指数按 $\epsilon$ 的幂次展开时，上式化为
\begin{equation}
  U(x) = 1 - ie^2\epsilon\bigl[ \partial_1A_2(x) - \partial_2A_1(x) \bigr]
    + \mathcal{O}(\epsilon^3) .
  \label{eq:15.12}
\end{equation}
因此结构
\begin{equation}
  F_{\mu\nu} = \partial_\mu A_\nu - \partial_\nu A_\mu
  \label{eq:15.13}
\end{equation}
是局域不变的。当然，$F_{\mu\nu}$ 就是熟悉的电磁场张量，它在 \eqref{eq:15.7} 下的不变性可以直接检验。然而，前面的构造向我们展示了 $F_{\mu\nu}$ 结构的几何起源。任何仅通过 $F_{\mu\nu}$ 及其导数依赖 $A_\mu$ 的函数都是局域不变的。更一般的函数，例如矢量场质量项 $m^2A_\mu A^\mu$，在 \eqref{eq:15.7} 下的变换方式无法被补偿，因而不能出现在不变拉格朗日量中。

利用协变导数可以对 $F_{\mu\nu}$ 的不变性给出一个相关的论证。我们上面已经看到，如果一个场具有局域变换律 \eqref{eq:15.1}，那么它的协变导数就具有相同的变换律。因此 $\psi$ 的二阶协变导数也按 \eqref{eq:15.1} 变换。对协变导数的对易子同样的结论也成立：
\begin{equation}
  [D_\mu, D_\nu]\psi(x) \to e^{i\alpha(x)}[D_\mu, D_\nu]\psi(x) .
  \label{eq:15.14}
\end{equation}
然而，这个对易子本身根本不是导数：
\begin{equation}
  \begin{aligned}
    [D_\mu, D_\nu]\psi &= \bigl[ \partial_\mu\partial_\nu\psi
      + ie(\partial_\mu A_\nu) \psi - ie A_\nu\partial_\mu\psi
      - ie A_\mu\partial_\nu\psi - e^2A_\mu A_\nu\psi \bigr] - (\mu\leftrightarrow\nu)\\
    &= ie\bigl( \partial_\mu A_\nu - \partial_\nu A_\mu \bigr)\cdot\psi .
  \end{aligned}
  \label{eq:15.16}
\end{equation}
也就是说，
\begin{equation}
  [D_\mu, D_\nu] = ieF_{\mu\nu} .
\end{equation}
在 \eqref{eq:15.14} 的右边，因子 $\psi(x)$ 承担了全部变换律，因此相乘的因子 $F_{\mu\nu}$ 必定不变。人们可以把协变导数的对易子想象为跨越一个小方块的比较之比较；因此，从根本上说，这个论证与上一段的论证是等价的。

无论用哪种方法证明 $F_{\mu\nu}$ 的不变性，我们现在已经集齐了为写出电子场 $\psi$ 及其相伴联络 $A_\mu$ 的最一般局域不变拉格朗日量所需的全部要素。这个拉格朗日量必须是 $\psi$ 及其协变导数、以及 $F_{\mu\nu}$ 及其导数的函数，并且必须对整体相位转动不变。一直到维数为 4 的算符，只可能有四项：
\begin{equation}
  \mathcal{L}_4 = \bar{\psi}(i\slashed{D})\psi - \frac{1}{4}F_{\mu\nu}F^{\mu\nu}
    - c\epsilon^{\mu\nu\rho\sigma}F_{\mu\nu}F_{\rho\sigma} - m\bar{\psi}\psi .
  \label{eq:15.17}
\end{equation}
通过调整场 $\psi$ 和 $A_\mu$ 的归一化，我们已经把前两项的系数设定为其标准值。对 $A_\mu$ 的这种归一化要求在我们对 $A_\mu$ 的原初定义 \eqref{eq:15.5} 中引入任意的标度因子 $e$。第三项破坏离散对称性 $P$ 和 $T$，因此如果我们假定这些对称性，就可以排除它。%
\footnote{$P$、$C$ 和 $T$ 破坏的一般系统学在第~20.3 节讨论。} 于是 $\mathcal{L}_4$ 只包含两个自由参数，即标度因子 $e$ 和系数 $m$。

通过使用维数为 5 和 6 的算符，我们可以构成许多额外的规范不变组合：
\begin{equation}
  \mathcal{L}_6 = i c_1\bar{\psi}\sigma^{\mu\nu}F_{\mu\nu}\psi
    + c_2\bigl( \bar{\psi}\psi \bigr)^2 + \cdots .
  \label{eq:15.18}
\end{equation}
在质量维数的每一更高阶还会出现更多允许的项。但所有这些项都是不可重整化的相互作用。用第~12.1 节的语言说，在截断趋于无穷的极限下，它们在四维中对物理是无关的。

\section{Yang-Mills 拉格朗日量}

在上一节中，我们从局域相位不变性的假设出发构建了 QED 拉格朗日量。这一构造中的基本要素是协变导数，它补偿了相邻点之间费米子场相位的变化。在本节中，我们将把这一构造推广到更广泛的一类局域对称性。

考虑一个在一个连续群的变换下具有整体对称性的理论。为明确起见，我们取这个群为 $SU(2)$，即行列式为 1 的 $2\times2$ 幺正矩阵构成的群。我们用一对复值 Dirac 场来表示理论的场：
\begin{equation}
  \psi(x) = \begin{pmatrix}\psi_1(x)\\ \psi_2(x)\end{pmatrix} .
\end{equation}
我们规定理论应在此对称性的局域形式下不变，
\begin{equation}
  \psi(x) \to V(x)\,\psi(x) ,
  \label{eq:15.20}
\end{equation}
其中 $V(x)$ 是行列式为 1 的 $2\times2$ 幺正矩阵，其元素随 $x$ 任意变化。和之前一样，导数 $\partial_\mu\psi$ 没有简单的变换律，因为相邻点处的场用不同的矩阵 $V$ 变换。

为了构造协变导数，我们再次引入比较器 $U(y,x)$，现在它必须是具有变换律
\begin{equation}
  U(y,x) \to V(y)\,U(y,x)\,V^{\dagger}(x) .
  \label{eq:15.22}
\end{equation}
的 $2\times2$ 矩阵。在无穷小间隔处，我们展开
\begin{equation}
  U(x+\epsilon n, x) = 1 - ig\epsilon n^\mu A_\mu(x) + \mathcal{O}(\epsilon^2),
  \label{eq:15.23}
\end{equation}
其中 $A_\mu(x)$ 现在是一个取 $2\times2$ 矩阵值的矢量场。我们把这样的量称为联络。于是协变导数为
\begin{equation}
  D_\mu\psi(x) = \bigl[ \partial_\mu + igA_\mu(x) \bigr]\psi(x) .
  \label{eq:15.24}
\end{equation}
由于 $V(x)$ 是幺正的，矩阵 $A_\mu(x)$ 可以取为 Hermitian 且无迹的。一个 Hermitian、无迹的 $2\times2$ 矩阵可以按 Pauli 矩阵基展开：
\begin{equation}
  A_\mu(x) = A_\mu^a(x)\ \frac{\sigma^a}{2} .
  \label{eq:15.25}
\end{equation}
这里指标 $a$ 从 1 取到 3，矩阵 $\sigma^a/2$ 是 $SU(2)$ 的生成元。因此，在 $SU(2)$ 规范理论中，有三个矢量场 $A_\mu^a$。

$A_\mu(x)$ 的变换律来自比较器的变换律。我们发现
\begin{equation}
  A_\mu(x)\ \frac{\sigma^a}{2} \to V(x)\Bigl( A_\mu(x)\frac{\sigma^a}{2}
    + \frac{1}{g}\partial_\mu \Bigr)\ V^{\dagger}(x) .
  \label{eq:15.27}
\end{equation}
导数作用于 $V^{\dagger}(x) = \exp(-i\alpha^a\sigma^a/2)$；要显式地计算这个导数并不那么容易，因为指数不一定与它的导数对易。对于无穷小变换，我们可以把 $V(x)$ 展开到 $\alpha$ 的一阶。在这种情况下我们得到
\begin{equation}
  A_\mu^a\ \frac{\sigma^a}{2} \to \Bigl( A_\mu^a + \frac{1}{g}\partial_\mu\alpha^a
    + \epsilon^{abc}\alpha^b A_\mu^c \Bigr)\ \frac{\sigma^a}{2} + \cdots .
  \label{eq:15.28}
\end{equation}
这个变换律中的最后一项是新的，它来自局域变换的不可对易性。把此关系与费米子变换的无穷小形式，
\begin{equation}
  \psi \to \Bigl( 1 + i\alpha^a\frac{\sigma^a}{2} \Bigr)\psi + \cdots ,
  \label{eq:15.29}
\end{equation}
相结合，我们可以检验协变导数的无穷小变换：
\begin{equation}
  D_\mu\psi \to \Bigl( 1 + i\alpha^a\frac{\sigma^a}{2} \Bigr)D_\mu\psi ,
  \label{eq:15.30}
\end{equation}
准确到 $\alpha^2$ 阶的项。利用 \eqref{eq:15.27} 和式~(15.21) 不难检验，即使对于有限变换，协变导数也与它所作用的场具有相同的变换律。

利用协变导数，我们可以构造涉及 $\psi$ 的最一般规范不变拉格朗日量。但要写出完整的拉格朗日量，我们还必须找到仅依赖 $A_\mu^i$ 的规范不变项。为此，我们构造电磁场张量的类比物。我们将使用上一节的第二种方法，即从协变导数的对易子出发。协变导数的变换律意味着
\begin{equation}
  [D_\mu, D_\nu]\psi(x) \to V(x)[D_\mu, D_\nu]\psi(x) .
  \label{eq:15.31}
\end{equation}
同时，利用公式 \eqref{eq:15.24} 写出对易子，我们可以像 Abelian 情形那样证明，$[D_\mu, D_\nu]$ 不是一个微分算符，而只是一个作用在 $\psi$ 上的相乘因子（现在是矩阵）。然而这一次有一个新特征：对易子展开式中的最后一项不再为零。相反，我们发现
\begin{equation}
  [D_\mu, D_\nu] = -igF^a_{\mu\nu}\ \frac{\sigma^a}{2} ,
  \label{eq:15.32}
\end{equation}
其中
\begin{equation}
  F^a_{\mu\nu} = \partial_\mu A_\nu^a - \partial_\nu A_\mu^a
    + g\epsilon^{abc}A_\mu^b A_\nu^c .
  \label{eq:15.33}
\end{equation}
我们可以应用 Pauli 矩阵的标准对易关系来简化这个关系：
\begin{equation}
  \Bigl[ \frac{\sigma^a}{2}, \frac{\sigma^b}{2} \Bigr] = i\epsilon^{abc}
    \frac{\sigma^c}{2} .
  \label{eq:15.34}
\end{equation}
于是
\begin{equation}
  F^a_{\mu\nu} = \partial_\mu A_\nu^a - \partial_\nu A_\mu^a
    + g\epsilon^{abc}A_\mu^b A_\nu^c .
  \label{eq:15.35}
\end{equation}
场强的变换律来自式~(15.21) 和 \eqref{eq:15.31}：
\begin{equation}
  F^a_{\mu\nu}\ \frac{\sigma^a}{2} \to V(x)\ F^b_{\mu\nu}\ \frac{\sigma^b}{2}\
    V^{\dagger}(x) .
  \label{eq:15.36}
\end{equation}
其无穷小形式为
\begin{equation}
  F^a_{\mu\nu} \to F^a_{\mu\nu} + i\alpha^b F^c_{\mu\nu}\ \epsilon^{bca} .
  \label{eq:15.37}
\end{equation}
注意场强不再是一个规范不变的量。它也不可能如此，因为现在有三个场强，每一个都与抽象空间中的某个给定转动方向相联系。然而，我们很容易构成场强的规范不变组合。例如，
\begin{equation}
  \frac{1}{4}F^a_{\mu\nu}F^{a\mu\nu}
  \label{eq:15.38}
\end{equation}
就是 $A_\mu^i$ 的一个规范不变动能项。注意，与电动力学的情形相反，这个拉格朗日量包含 $A_\mu^i$ 的三次和四次项。因此，这个拉格朗日量描述一个非平凡的相互作用场论，称为\emph{Yang-Mills 理论}。这是非阿贝尔规范理论最简单的例子。

为了构造与费米子相互作用的 Yang-Mills 矢量场的理论，我们只需把规范场拉格朗日量 \eqref{eq:15.38} 加到熟悉的 Dirac 拉格朗日量上，并把 $\psi$ 的普通导数替换为协变导数。结果看起来几乎与 QED 拉格朗日量相同：
\begin{equation}
  \mathcal{L} = \bar{\psi}(i\slashed{D} - m)\psi - \frac{1}{4}(F^a_{\mu\nu})^2 .
  \label{eq:15.39}
\end{equation}
这就是著名的 Yang-Mills 拉格朗日量。和 QED 的拉格朗日量一样，它依赖两个参数：标度因子 $g$（类似于电子电荷）和费米子质量 $m$。对这个拉格朗日量变分，我们得到规范理论的经典运动方程。它们是对费米子场的 Dirac 方程和对矢量场的方程
\begin{equation}
  \partial^\mu F^a_{\mu\nu} + g\epsilon^{abc}A^{b\mu}F^c_{\mu\nu} = gj^a_\nu
  \label{eq:15.40}
\end{equation}

我们对 $SU(2)$ 对称性变换 \eqref{eq:15.20} 所做的一切都很容易推广到任何其他的连续对称性群。可能的对称群的完整范围将在第~15.4 节中枚举和分类。然而，对任何这样的群，拉格朗日量各要素的一般表达式都非常相似。考虑任何一个由一组 $n\times n$ 幺正矩阵 $V$ 表示的连续变换群。于是基本场 $\psi(x)$ 构成一个 $n$ 重态，并按如下方式变换：
\begin{equation}
  \psi(x) \to V(x)\psi(x) ,
  \label{eq:15.41}
\end{equation}
其中 $V$ 对 $x$ 的依赖使得变换是局域的。在无穷小形式中，$V(x)$ 可以用对称性群的一组基本生成元展开，这些生成元可以表示为 Hermitian 矩阵 $t^a$：
\begin{equation}
  V(x) = 1 + i\alpha^a(x)\,t^a + \mathcal{O}(\alpha^2) .
  \label{eq:15.42}
\end{equation}
现在，对于一般的局域对称性群，我们只需在分析的每一步用
\begin{equation}
  \frac{\sigma^a}{2} \to t^a
  \label{eq:15.43}
\end{equation}
作替换，就可以把从式~\eqref{eq:15.22} 到式~\eqref{eq:15.33} 的整个分析贯彻到底。

为了推广场张量的显式表达式 \eqref{eq:15.35}，我们需要知道矩阵 $t^a$ 的对易关系。通常把这些对易关系写成标准形式
\begin{equation}
  [t^a, t^b] = if^{abc}t^c ,
  \label{eq:15.44}
\end{equation}
其中 $f^{abc}$ 是一组数，称为\emph{结构常数}。这个对象在式~\eqref{eq:15.34} 中取代 $\epsilon^{ijk}$。通常要选择矩阵 $t^a$ 的一个基，使得 $f^{abc}$ 完全反对称；我们将在第~15.4 节证明这总是可能的。

我们现在可以如下重述我们所有的结果。与一般变换 \eqref{eq:15.41} 相联系的协变导数为
\begin{equation}
  D_\mu = \partial_\mu - igA_\mu^a t^a ;
  \label{eq:15.45}
\end{equation}
它对局域对称性的每一个独立生成元都含有一个矢量场。$\psi$ 和 $A_\mu^a$ 的无穷小变换律为
\begin{equation}
  \psi \to \bigl( 1 + i\alpha^a t^a \bigr)\psi;\qquad
  A_\mu^a \to A_\mu^a + \frac{1}{g}\partial_\mu\alpha^a
    + f^{abc}\alpha^b A_\mu^c .
  \label{eq:15.46}
\end{equation}
$A_\mu^a$ 的有限变换恰具有 \eqref{eq:15.27} 的形式：
\begin{equation}
  A_\mu^a(x)t^a \to V(x)\Bigl( A_\mu^a(x)t^a + \frac{1}{g}\partial_\mu
    \Bigr)\ V^{\dagger}(x) .
  \label{eq:15.47}
\end{equation}
这些变换律意味着 $\psi$ 的协变导数与 $\psi$ 本身具有相同的变换律。场张量定义为
\begin{equation}
  [D_\mu, D_\nu] = -igF^a_{\mu\nu}t^a ,
  \label{eq:15.48}
\end{equation}
或更显式地，
\begin{equation}
  F^a_{\mu\nu} = \partial_\mu A_\nu^a - \partial_\nu A_\mu^a
    + gf^{abc}A_\mu^b A_\nu^c .
  \label{eq:15.49}
\end{equation}
这个量具有无穷小变换
\begin{equation}
  F^a_{\mu\nu} \to F^a_{\mu\nu} + f^{abc}\alpha^b F^c_{\mu\nu} .
  \label{eq:15.50}
\end{equation}
从式~\eqref{eq:15.46} 和 \eqref{eq:15.50} 可以证明，$\psi$、$D_\mu\psi$ 及其协变导数的任何整体对称函数也是局域对称的，因此是规范不变拉格朗日量中一个项的候选者。然而，一直到维数 4，允许的项非常少。最一般的、可重整化且守恒 $P$ 和 $T$ 的规范不变拉格朗日量仍由式~\eqref{eq:15.39} 给出。相应的经典运动方程为
\begin{equation}
  \partial^\mu F^a_{\mu\nu} + gf^{abc}A^{b\mu}F^c_{\mu\nu} = -gj^a_\nu ,
  \label{eq:15.51}
\end{equation}
其中
\begin{equation}
  j^{a\mu} = \bar{\psi}\gamma^\mu t^a\psi
  \label{eq:15.52}
\end{equation}
是费米子场的整体对称性流。

注意 Yang-Mills 拉格朗日量 \eqref{eq:15.39} 中的非线性项出现在协变导数中（在那里它们正比于 $t^a$）以及场张量中（在那里它们正比于 $f^{abc}$）。因此，非阿贝尔规范理论中相互作用的形式由局域对称性所决定。矢量场与自身的非线性相互作用正比于对称性生成元的对易子，从而明确地要求对称性群具有非阿贝尔性质。

\section{规范不变的 Wilson 圈}

在前两节中我们都用到了比较器 $U(y,x)$，它把 $x$ 点处费米子的规范变换律转换为 $y$ 点处的规范变换律。到目前为止，在写出这个对象的表达式时，假定 $x$ 和 $y$ 为无穷小间隔就足够了。然而，值得进一步考虑 $x$ 和 $y$ 相距很远时比较器的情况。这一讨论将使我们更深入地理解规范不变性的几何，并揭示规范场的一些额外有用的函数，我们将在第~19 章中运用它们。

我们首先回到 Abelian 理论，并在该情境下展开我们对 $U(y,x)$ 的讨论。在式~\eqref{eq:15.10} 中，我们构造了一个沿围绕小方块的路径上的比较器乘积。我们证明了即使最终回到出发点，这个乘积 $U(x)$ 也并不是平凡的；相反，我们发现 $U(x)$ 与 1 相差一个正比于电磁场强和小方块面积的项。这是一般结论的一个特例：间隔有限的两点 $x$ 和 $y$ 之间的比较器依赖于从 $x$ 到 $y$ 所取的路径。

为了解释这一论断，反转第~15.1 节的部分逻辑是有用的。我们从联络 $A_\mu$ 出发，假设它具有变换律 \eqref{eq:15.7}，并构造 $U(z,y)$ 作为 $A_\mu$ 的按 \eqref{eq:15.3} 变换的函数。不难验证，如果积分沿任何从 $y$ 到 $z$ 的路径 $P$ 进行，则表达式
\begin{equation}
  U_P(z,y) = \exp\Bigl[ -ie\int_y^z dx^\mu\ A_\mu(x) \Bigr]
  \label{eq:15.53}
\end{equation}
满足这一判据。这个对象 $U_P(z,y)$ 称为\emph{Wilson 线}。%
\footnote{这个依赖路径的相位早在 Wilson 的工作之前就已使用，见于 Schwinger 关于 QED 的早期论文，以及 Y.~Aharonov 和 D.~B\"ohm, \emph{Phys.\ Rev.} \textbf{115}, 485 (1959)。} 表达式 \eqref{eq:15.53} 为有限间隔的点给出了抽象比较器 $U(z,y)$ 的一个显式实现。

Wilson 线的一个关键性质是它依赖于路径 $P$。如果 $P$ 是回到 $y$ 的闭合路径，我们就得到\emph{Wilson 圈}，
\begin{equation}
  U_P(y,y) = \exp\Bigl[ -ie\oint_P dx^\mu\ A_\mu(x) \Bigr] .
  \label{eq:15.54}
\end{equation}
这个量是 $A_\mu$ 的一个非平凡函数，并且按构造是局域规范不变的。事实上，$A_\mu$ 的所有规范不变函数都可以看作对路径 $P$ 的各种选取的 Wilson 圈之组合。为了论证这一论断，我们用 Stokes 定理把 Wilson 圈改写为
\begin{equation}
  \exp\Bigl[ -ie\oint_P dx^\mu\ A_\mu(x) \Bigr]
  = \exp\Bigl[ -\frac{ie}{2}\int_\Sigma d\sigma^{\mu\nu}\ F_{\mu\nu} \Bigr],
  \label{eq:15.55}
\end{equation}
其中 $\Sigma$ 是张在闭合路径 $P$ 上的曲面，
$d\sigma^{\mu\nu}$ 是这个曲面上的面积元，$F$ 是场张量 \eqref{eq:15.13}。Wilson 圈与场强之间的这种关系如图~\ref{fig:15.2} 所示。由于 Wilson 圈是规范不变的，这个论证给出可视化场强规范不变性的又一种方式。反过来，由于 $A_\mu$ 的（几乎）所有规范不变函数都可以由 $F_{\mu\nu}$ 构造出来，这个表达式也支持 $U_P(y,y)$ 是最一般的规范不变量的论断。

\begin{figure}[t]
  \centering
  \includegraphics[width=0.5\textwidth]{fig15.2.png}
  \caption{Wilson 圈与场强之间的关系。}
  \label{fig:15.2}
\end{figure}

Wilson 线和 Wilson 圈都可以推广到非阿贝尔情形。然而，在这里，当我们考虑不可对易矩阵的指数时，会出现额外的微妙之处。让我们首先构造 Wilson 线，它现在按式~\eqref{eq:15.22} 变换。把 \eqref{eq:15.53} 直接改写成指数中含 $A_\mu^a t^a$ 的积分是不正确的，因为这些矩阵在不同点处不一定对易。相反，我们必须以特定的方式给这些矩阵编序。我们现在就给出正确的编序规则，然后证明它的变换律。

设 $s$ 是路径 $P$ 的一个参数，在 $x=y$ 处从 0 取到 $x=z$ 处的 1。于是把 Wilson 线定义为指数的幂级数展开，其中每一项中的矩阵都按 $s$ 值较大的排在左边的规则编序：
\begin{equation}
  U_P(z,y) = 1 + (-ig)\int_0^1 ds\ \frac{dx^\mu}{ds}\ A_\mu(x(s))
    + (-ig)^2\int_0^1 ds_1\int_0^{s_1} ds_2\ \frac{dx^\mu}{ds_1}\ A_\mu(x(s_1))
      \ \frac{dx^\nu}{ds_2}\ A_\nu(x(s_2)) + \cdots .
  \label{eq:15.56}
\end{equation}
这个对象称为\emph{路径编序指数}，通常写成
\begin{equation}
  U_P(z,y) = P\ \exp\Bigl[ -ig\int_y^z dx^\mu\ A_\mu(x) \Bigr] ,
  \label{eq:15.57}
\end{equation}
其中 $P$（代表「路径」）指示我们沿路径给矩阵编序。有了这个定义，可以验证 $U_P(z,y)$ 按 \eqref{eq:15.22} 变换。于是非阿贝尔 Wilson 圈为
\begin{equation}
  U_P(y,y) = P\ \exp\Bigl[ -ig\oint_P dx^\mu\ A_\mu(x) \Bigr] .
  \label{eq:15.58}
\end{equation}
这个对象的迹，
\begin{equation}
  \mathrm{tr}\Bigl[ P\ \exp\Bigl( -ig\oint dx^\mu\ A_\mu \Bigr) \Bigr],
\end{equation}
是规范场的一个规范不变函数，因为迹在循环置换下不变。这个量在非阿贝尔规范理论的非微扰研究中起着核心作用。

\section{关于李代数的一些基本事实}

在第~15.2 节的末尾我们看到，非阿贝尔规范理论的类别非常庞大。为了最有效地处理这些理论，值得暂停一下，考虑它们所基于的连续群的一般性质。在本节中，我们将枚举所有可以用来构造非阿贝尔规范理论的可能的群。然后我们将计算一些由群变换矩阵构造出来的数值因子，它们在进行量子化规范理论的显式计算时是需要的。%
\footnote{在本节中，我们将不加证明地陈述连续群理论的一些一般结果。有许多优秀的著作系统地综述了这些数学结果。其中，我们特别推荐 Cahn (1984)，它给出简短而精辟的讨论；以及 S.~Helgason, \emph{Differential Geometry, Lie Groups, and Symmetric Spaces} (Academic Press, 1978)，它给出优雅而严谨的论述。R.~Slansky, \emph{Phys.\ Repts.} \textbf{79}, 1 (1981) 编纂了一组特别有用的与构造非阿贝尔规范理论相关的群论恒等式表。}

对数学家而言，群由遵守一定代数规则的抽象实体组成。然而，在量子力学中，我们特别感兴趣的是作用于量子态矢量空间上的幺正算符群。我们关注连续生成的群，即包含任意接近恒等元的元素、使得一般元素可以通过这些无穷小元素的反复作用而达到的群。于是任意无穷小群元素 $g$ 都可以写成
\begin{equation}
  g(\alpha) = 1 + i\alpha^a T^a + \mathcal{O}(\alpha^2) .
  \label{eq:15.67}
\end{equation}
无穷小群参数 $\alpha^a$ 的系数是 Hermitian 算符 $T^a$，称为对称性群的\emph{生成元}。具有这种结构的连续群称为\emph{李群}。

生成元的集合 $T^a$ 必须张成无穷小群变换的空间，因此生成元 $T^a$ 的对易子必须是生成元的线性组合。于是算符 $T^a$ 的对易关系可以写成
\begin{equation}
  [T^a, T^b] = if^{abc}T^c ;
  \label{eq:15.68}
\end{equation}
数 $f^{abc}$ 称为\emph{结构常数}。由生成元张成、并带有额外的对易运算的矢量空间称为\emph{李代数}。

对易关系 \eqref{eq:15.68} 和恒等式
\begin{equation}
  [T^a, [T^b, T^c]] + [T^b, [T^c, T^a]] + [T^c, [T^a, T^b]] = 0 ,
  \label{eq:15.70}
\end{equation}
称为\emph{Jacobi 恒等式}，它们定义了李代数。从数学家的观点看（把生成元看作抽象实体而不是 Hermitian 算符），Jacobi 恒等式是一个公理，一组给定的对易规则必须满足它才能定义李代数。

李代数的对易关系完全确定了相伴李群在充分接近恒等元处的群乘法律。对于足够大的变换，额外的整体性问题会起作用；举一个熟悉的例子，$SU(2)$ 和 $O(3)$ 具有相同的对易关系但不同的整体结构。然而，非阿贝尔规范理论的拉格朗日量只依赖于局域对称性群的李代数，因此我们从现在起将忽略这些整体性问题。

\subsection{李代数的分类}

对于规范理论的应用，局域对称性通常是一组场的幺正变换。因此我们主要关心具有有限维 Hermitian 表示的李代数，它们导致相应李群的有限维幺正表示。我们还将假设生成元的数目有限。这样的李代数称为\emph{紧致}的，因为这些条件意味着李群是一个有限维紧致流形。

如果生成元 $T^a$ 中有一个与所有其他生成元对易，它就生成一个独立的连续 Abelian 群。这样的群具有相位转动群
\begin{equation}
  \psi \to e^{i\alpha}\psi ,
  \label{eq:15.71}
\end{equation}
的结构，我们称之为 $U(1)$。如果代数中不含这样的可对易元素，即群中不含 $U(1)$ 因子，那么我们称该代数为\emph{半单}的。此外，如果李代数不能分成两组互相对易的生成元，该代数就是\emph{单}的。一般李代数是不可约非阿贝尔单分量与额外 Abelian 生成元的直和。

令人惊讶的是，李代数为紧致且单的这些基本条件竟然极其苛刻。作为十九世纪数学的一项成就，Killing 和 Cartan 对所有可能的紧致单李代数作了分类。这些代数几乎全部属于三个无穷族，只有五个例外。这三个无穷族是所谓\emph{经典群}对应的代数，其结构可以方便地用特定的矩阵表示来定义。三个经典群族的定义如下：
\begin{enumerate}
  \item \emph{$N$ 维矢量的幺正变换。} 设
    $\xi$ 和 $\eta$ 是复 $N$ 维矢量。于是一般线性变换具有形式
    \begin{equation}
      \eta_a \to U_{ab}\eta_b,\qquad \xi_a \to U_{ab}\xi_b .
      \label{eq:15.72}
    \end{equation}
    如果该变换保持内积
    $\eta^{*}_a\xi_a$，我们就说它是幺正的。纯相位变换
    \begin{equation}
      \xi_a \to e^{i\alpha}\xi_a
      \label{eq:15.73}
    \end{equation}
    构成一个与所有其他幺正变换对易的 $U(1)$ 子群；我们去掉这个子群以构成一个单李群，
    称为 $SU(N)$；它由所有满足 $\det(U) = 1$ 的 $N\times N$ 幺正变换组成。$SU(N)$ 的生成元用
    $N\times N$ Hermitian 矩阵 $t^a$ 表示，条件是它们与 \eqref{eq:15.73} 的生成元正交：
    \begin{equation}
      \mathrm{tr}[t^a] = 0 .
      \label{eq:15.74}
    \end{equation}
    满足这些条件的独立矩阵有 $N^2-1$ 个。
  \item \emph{$N$ 维矢量的正交变换。} 这是
    保持对称内积
    \begin{equation}
      \eta_a E_{ab}\xi_b,
      \qquad\text{with } E_{ab} = \delta_{ab} .
      \label{eq:15.75}
    \end{equation}
    的幺正 $N\times N$ 变换子群。
    这就是通常的矢量内积，因此这个群是 $N$ 维空间中的转动群
    $SO(N)$。（加入反射就得到群
    $O(N)$。）$N$ 维空间中的每一个平面都对应一个独立的转动，所以 $SO(N)$ 有 $N(N-1)/2$ 个生成元。
  \item \emph{$N$ 维矢量的辛变换。} 这是
    对 $N$ 为偶数、保持反对称内积
    \begin{equation}
      \eta_a E_{ab}\xi_b,
      \qquad\text{with } E_{ab} = \begin{pmatrix}0&1\\-1&0\end{pmatrix}
      \text{ (in } N/2\times N/2 \text{ blocks)},
      \label{eq:15.76}
    \end{equation}
    的幺正 $N\times N$ 变换子群，
    其中矩阵的元素是 $N/2\times N/2$ 的块。这个群
    称为 $Sp(N)$；它有 $N(N+1)/2$ 个生成元。
\end{enumerate}

在这三个族之外，还有五个\emph{例外}李代数，在 Cartan 的分类体系中记为 $G_2$、$F_4$、$E_6$、$E_7$ 和 $E_8$。其中，$E_6$ 和 $E_8$ 已被用作基本相互作用的有趣统一模型中的局域对称性群。然而，在本书中我们将不再进一步考虑这些例外群。事实上，我们的大多数例子将只涉及 $SU(N)$ 群。

\subsection{表示}

一旦我们指定了局域对称性群，出现在拉格朗日量中的场最自然地按这个群的一个有限维幺正表示变换。因此，我们接下来可能问，如何系统地找到任何给定李群的所有这样的表示。回想对于群 $SU(2)$，表示可以直接从对易关系出发，利用升算符和降算符 $J^+$ 和 $J^-$ 构造。这个构造可以推广以找到任何紧致李代数的有限维表示。然而，在本书中，我们将处理相对简单的表示，其结构可以用不太形式化的方法推得。

在讨论李代数的表示之前，我们应当回顾群表示的一些一般方面。给定对称性群 $G$，该群的李代数的一个有限维幺正表示是一组满足对易关系 \eqref{eq:15.68} 的 $d\times d$ Hermitian 矩阵 $t^a$。大小 $d$ 是表示的\emph{维数}。任意表示通常可以通过找到一个使所有表示矩阵同时为分块对角的基来分解。通过这个基变换，我们可以把表示写成\emph{不可约表示}的直和。我们用 $t^a_r$ 记不可约表示 $r$ 中的表示矩阵。

对矩阵 $t^a$ 采用基于其乘积之迹的归一化约定是标准做法。如果李代数是半单的，矩阵 $t^a_r$ 本身就是无迹的。然而，考虑两个生成元矩阵之积的迹：
\begin{equation}
  \mathrm{tr}\bigl[ t^a_r t^b_r \bigr] = D^{ab} .
  \label{eq:15.77}
\end{equation}
只要生成元矩阵是 Hermitian 的，矩阵 $D^{ab}$ 就是正定的。让我们为生成元 $T^a$ 选择一个基，使这个矩阵正比于恒等矩阵。可以证明，一旦对某一个不可约表示做到这一点，它对所有不可约表示都成立。因此，在这个基下，
\begin{equation}
  \mathrm{tr}\bigl[ t^a_r t^b_r \bigr] = C(r)\,\delta^{ab} ,
  \label{eq:15.78}
\end{equation}
其中 $C(r)$ 对每个表示 $r$ 是一个常数。式
\eqref{eq:15.78} 和对易关系 \eqref{eq:15.68} 给出结构常数的如下表示：
\begin{equation}
  f^{abc} = -\frac{i}{C(r)}\,\mathrm{tr}\bigl( \bigl[ t^a_r, t^b_r \bigr]
    t^c_r \bigr) .
  \label{eq:15.79}
\end{equation}
这个式子意味着 $f^{abc}$ 是完全反对称的。

对 $G$ 的每一个不可约表示 $r$，都有一个相联系的\emph{共轭表示} $\bar{r}$。表示 $r$ 给出无穷小变换
\begin{equation}
  \phi \to \bigl( 1 + i\alpha^a t^a_r \bigr)\phi .
  \label{eq:15.80}
\end{equation}
这个变换的复共轭，
\begin{equation}
  \phi^{*} \to \bigl( 1 - i\alpha^a (t^a_r)^{*} \bigr)\phi^{*} ,
  \label{eq:15.81}
\end{equation}
也必须是 $G$ 的一个表示的无穷小元素。因此 $r$ 的共轭表示具有表示矩阵
\begin{equation}
  t^a_{\bar{r}} = -(t^a_r)^{*} = -(t^a_r)^T .
  \label{eq:15.82}
\end{equation}
由于 $\phi^{*}\phi$ 在幺正变换下不变，可以把分别在表示 $r$ 和 $\bar{r}$ 中变换的场组合起来构成一个群不变量。

如果存在幺正变换 $U$ 使得
$t^a_{\bar{r}} = Ut^a_rU^{\dagger}$，那么表示 $\bar{r}$ 可能等价于 $r$。
如果是这样，表示 $r$ 就是\emph{实的}。在这种情况下，存在一个矩阵 $G_{ab}$，使得如果 $\eta$ 和
$\xi$ 属于表示 $r$，则组合
$G_{ab}\eta_a\xi_b$ 是一个不变量。区分 $G_{ab}$ 对称和 $G_{ab}$ 反对称这两种情形有时是有用的。在前一种情形下，表示是\emph{严格实}的；在后一种情形下它是\emph{伪实}的。这两种情形在 $SU(2)$ 中都已经出现：两个矢量的不变组合是 $v^aw^a$，所以矢量是实表示；两个旋量的不变组合是 $\epsilon^{ab}\chi_a\psi_b$，所以旋量是伪实表示。

用这些语言，我们可以讨论经典群的最简单表示。在 $SU(N)$ 中，基本不可约表示（通常称为\emph{基本表示}）是 $N$ 维复矢量。对于 $N > 2$，这个表示是复的，因此还有第二个不等价的表示 $\bar{N}$。（在 $SU(2)$ 中，这个表示就是伪实旋量表示。）在 $SO(N)$ 中，基本的 $N$ 维矢量是（严格）实表示。在 $Sp(N)$ 中，$N$ 维矢量是伪实表示。

另一个对任何单李代数都存在的不可约表示，是代数的生成元所属的表示。这个表示称为\emph{伴随表示}，记为 $r = G$。其表示矩阵由结构常数给出：
\begin{equation}
  (t^b_G)_{ac} = if^{abc} .
  \label{eq:15.83}
\end{equation}
有了这个定义，$t^b_G$ 满足李代数
\begin{equation}
  \bigl( [t^b_G, t^c_G] \bigr)_{ae} = if^{bcd}(t^d_G)_{ae}
  \label{eq:15.84}
\end{equation}
的说法就只是 Jacobi 恒等式 \eqref{eq:15.70} 的一种改写。由于结构常数是实的且反对称的，
$t^a_G = -(t^a_G)^{*}$；因此伴随表示总是实表示。由上面给出的李群描述，伴随表示的维数 $d(G)$ 对经典群由下式给出：
\begin{equation}
  d(G) = \begin{cases}
    N^2-1 & \text{for } SU(N),\\
    N(N-1)/2 & \text{for } SO(N),\\
    N(N+1)/2 & \text{for } Sp(N) .
  \end{cases}
  \label{eq:15.85}
\end{equation}

把 $f^{abc}$ 视为表示矩阵，使我们能够更深入地理解第~15.2 节中引入的一些量。作用于伴随表示中场的协变导数为
\begin{equation}
  \begin{aligned}
    (D_\mu\phi)^a &= \partial_\mu\phi^a - igA_\mu^b(t^b_G)_{ac}\phi^c\\
    &= \partial_\mu\phi^a + gf^{abc}A_\mu^b\phi^c .
  \end{aligned}
  \label{eq:15.86}
\end{equation}
于是我们可以把 \eqref{eq:15.46} 中矢量场规范变换的无穷小形式识别为场在伴随表示下的运动。

另一个重要的群论量是\emph{二次 Casimir 算符}。对给定的不可约表示 $r$，我们定义
\begin{equation}
  C_2(r) = \sum_a t^a_r\,t^a_r .
  \label{eq:15.92}
\end{equation}
根据 Schur 引理，这个矩阵与所有生成元 $t^a_r$ 对易，因此对于不可约表示，它必须正比于恒等矩阵：
\begin{equation}
  C_2(r) = C_2(r)\cdot\mathbb{1} .
  \label{eq:15.93}
\end{equation}
比例常数 $C_2(r)$ 是 Casimir 算符在表示 $r$ 中的本征值。对 \eqref{eq:15.92} 取迹并利用 \eqref{eq:15.78}，我们发现
\begin{equation}
  C_2(r)\cdot d(r) = C(r)\cdot d(G) ,
  \label{eq:15.94}
\end{equation}
其中 $d(r)$ 是表示 $r$ 的维数。对于伴随表示，Casimir 算符为
\begin{equation}
  C_2(G) = C(G) .
\end{equation}
对于 $SU(N)$ 的基本表示，可以直接计算 $C_2(r)$。$SU(N)$ 的基本表示作用于 $N$ 维复矢量，其生成元按一般的无迹 $N\times N$ 张量变换；实现这些变换的矩阵构成了 $SU(N)$ 的伴随表示。在这种情况下，基本表示与共轭表示之积的分解给出
\begin{equation}
  N\times\bar{N} = 1 + (N^2-1) .
  \label{eq:15.102}
\end{equation}
对于这个分解，式~(15.100) 和式~(15.101) 意味着恒等式
\begin{equation}
  \Bigl( 2C_2(r)\cdot N \Bigr) \cdot N^2 = 0 + C_2(G)\cdot(N^2-1) .
  \label{eq:15.103}
\end{equation}
因此，对 $SU(N)$，
\begin{equation}
  C_2(G) = C(G) = N .
  \label{eq:15.104}
\end{equation}
二次 Casimir 算符计算的其他一些例子在习题~15.5 中给出。然而，我们在本节中讨论的例子，加上我们所回顾的基本群论概念，已经提供了足够的材料来执行非阿贝尔规范理论中最有物理意义的计算。

\problems

\textbf{15.1} \emph{$SU(3)$ 中的蛮力计算。}  $SU(3)$ 基本表示的标准基为
\begin{equation}
  \lambda^1 = \begin{pmatrix}0&1&0\\1&0&0\\0&0&0\end{pmatrix},\qquad
  \lambda^2 = \begin{pmatrix}0&-i&0\\i&0&0\\0&0&0\end{pmatrix},\qquad
  \lambda^3 = \begin{pmatrix}1&0&0\\0&-1&0\\0&0&0\end{pmatrix},\qquad
  \cdots
\end{equation}
\begin{enumerate}[label=(\alph*)]
  \item 解释为什么这个基中恰好有八个矩阵。
  \item 求所有这些矩阵的对易子，以确定 $SU(3)$ 的结构常数。证明，用这里的归一化，$f^{abc}$ 是完全反对称的。（这个练习很繁琐；你可能只想检验其中一组有代表性的对易子。）
  \item 检验正交性条件 \eqref{eq:15.78}，并求这个表示的常数 $C(r)$。
  \item 直接从二次 Casimir 算符 $C_2(r)$ 的定义 \eqref{eq:15.92} 计算它，并验证 $C_2(r)$ 与 $C(r)$ 之间的关系 \eqref{eq:15.94}。
\end{enumerate}

\textbf{15.2} 写出 $SU(2)$ 伴随表示的基矩阵。直接从它们的定义 \eqref{eq:15.78} 和 \eqref{eq:15.92} 计算 $C(G)$ 和 $C_2(G)$。

\textbf{15.3} \emph{库仑势。}
\begin{enumerate}[label=(\alph*)]
  \item 利用泛函积分，计算不含费米子的纯量子电动力学中 Wilson 圈的期望值。证明
    \begin{equation}
      \langle U_P(z,z) \rangle = \exp\Bigl[ -\frac{e^2}{2}\oint_P dx^\mu\;
        \oint_P dy^\nu\ g_{\mu\nu}\ \frac{1}{8\pi^2(x-y)^2} \Bigr] ,
    \end{equation}
    其中 $x$ 和 $y$ 沿闭合曲线 $P$ 积分。
  \item 考虑一个（类空的）宽度为 $R$、（类时的）长度为 $T$、$T\gg R$ 的矩形路径的 Wilson 圈。在这个极限下计算 Wilson 圈的期望值，并把它与时间演化的普遍表达式
    \begin{equation}
      \langle U_P \rangle = \exp[-iE(R)T] ,
    \end{equation}
    相比较，其中 $E(R)$ 是对应于 Wilson 圈的电磁源的能量。证明这些源的势能正好是库仑势，$V(R) = -e^2/4\pi R$。
  \item 假设非阿贝尔规范场的传播子由 Feynman 规范表达式
    \begin{equation}
      \frac{-i\delta^{ab}g_{\mu\nu}}{q^2+i\epsilon} ,
    \end{equation}
    给出，计算非阿贝尔规范理论中一个小矩形 Wilson 圈的领头阶期望值，并证明结果正比于源的表示之 Casimir 算符 $C_2(r)$。
\end{enumerate}
